% GENERATED FILE. Edit canonical lessons, then run npm run generate:books.
% canonical-source-hash: fnv1a64:cc20aa00e2382c75
% canonical-lessons: ES-C699-rayo, ES-C699-relampaguear, ES-C699-granizo, ES-C699-granizar, ES-C699-tronar

\chapter{Rayos y truenos}
\label{ch:es-rayos-y-truenos}
\hlchaptermodality{699}

\begin{chapteropening}
\textbf{By the end of this chapter:} \emph{I can talk about a storm: lightning, thunder and hail.}

Complete the last lesson of chapter 699: I can talk about a storm: lightning, thunder and hail.
\end{chapteropening}

\section[el rayo]{el rayo — the lightning bolt; the ray}

\label{lesson:ES-C699-rayo}



\begin{quote}
Before the new one: say the Spanish for to cloud over, then the Spanish for to clear.
\end{quote}

\subsection*{You'll want to know: el rayo}
\textbf{el rayo} — \textquotedblleft{}the lightning bolt; the ray\textquotedblright{}.

\begin{quote}
Un rayo cayó en el árbol.
\end{quote}

\begin{quote}
\emph{Lightning struck the tree.}
\end{quote}

\emph{Los rayos del sol} — the sun’s rays.

\begin{cousinweb}
\emph{rayo} is inherited from Latin \emph{radius} \textquotedblleft{}a staff, a spoke, a ray\textquotedblright{}.
\end{cousinweb}

\begin{grammarlens}[title={What you've built}]
One more word for telling, explaining and describing at B1.
\end{grammarlens}

\subsection*{Your turn}
\begin{itemize}
\raggedright
  \item \emph{Say it:} \emph{el rayo}
  \item \emph{Say it:} \emph{el rayo}, once more
  \item \emph{Say it:} a sentence of your own with \emph{el rayo}
  \item \emph{Recall:} say the Spanish for to cloud over, then the Spanish for to clear, then say \emph{el rayo} again
\end{itemize}

\begin{tcolorbox}[breakable,colback=teal!4,colframe=teal!35!black,title={Before you move on}]
What does el rayo mean? (\textquotedblleft{}The lightning bolt; the ray\textquotedblright{}.) Say it once more.
\end{tcolorbox}
\section[relampaguear]{relampaguear — to flash with lightning}

\label{lesson:ES-C699-relampaguear}



\begin{quote}
Before the new one: say the Spanish for to clear, then the Spanish for the lightning bolt.
\end{quote}

\subsection*{You'll want to know: relampaguear}
\textbf{relampaguear} — \textquotedblleft{}to flash with lightning\textquotedblright{}.

\begin{quote}
Toda la noche estuvo relampagueando.
\end{quote}

\begin{quote}
\emph{There was lightning all night.}
\end{quote}

Used in the third person: \emph{relampaguea}.

\begin{cousinweb}
\emph{relampaguear} is built on \emph{relámpago} \textquotedblleft{}flash of lightning\textquotedblright{} (EtymDB), from Latin \emph{lampas} \textquotedblleft{}a torch\textquotedblright{}, Greek \emph{lampás}.
\end{cousinweb}

\begin{grammarlens}[title={What you've built}]
One more word for telling, explaining and describing at B1.
\end{grammarlens}

\subsection*{Your turn}
\begin{itemize}
\raggedright
  \item \emph{Say it:} \emph{relampaguear}
  \item \emph{Say it:} \emph{relampaguear}, once more
  \item \emph{Say it:} a sentence of your own with \emph{relampaguear}
  \item \emph{Recall:} say the Spanish for to clear, then the Spanish for the lightning bolt, then say \emph{relampaguear} again
\end{itemize}

\begin{tcolorbox}[breakable,colback=teal!4,colframe=teal!35!black,title={Before you move on}]
What does relampaguear mean? (\textquotedblleft{}To flash with lightning\textquotedblright{}.) Say it once more.
\end{tcolorbox}
\section[el granizo]{el granizo — the hail}

\label{lesson:ES-C699-granizo}



\begin{quote}
Before the new one: say the Spanish for the lightning bolt, then the Spanish for to flash with lightning.
\end{quote}

\subsection*{You'll want to know: el granizo}
\textbf{el granizo} — \textquotedblleft{}the hail\textquotedblright{}.

\begin{quote}
El granizo rompió los cristales.
\end{quote}

\begin{quote}
\emph{The hail broke the windows.}
\end{quote}

One hailstone is \emph{un grano de granizo} or \emph{una piedra}.

\begin{cousinweb}
\emph{granizo} is built on \emph{grano} \textquotedblleft{}grain\textquotedblright{}, which is inherited from Latin \emph{grānum} (EtymDB): hail falls like grains of ice.
\end{cousinweb}

\begin{grammarlens}[title={What you've built}]
One more word for telling, explaining and describing at B1.
\end{grammarlens}

\subsection*{Your turn}
\begin{itemize}
\raggedright
  \item \emph{Say it:} \emph{el granizo}
  \item \emph{Say it:} \emph{el granizo}, once more
  \item \emph{Say it:} a sentence of your own with \emph{el granizo}
  \item \emph{Recall:} say the Spanish for the lightning bolt, then the Spanish for to flash with lightning, then say \emph{el granizo} again
\end{itemize}

\begin{tcolorbox}[breakable,colback=teal!4,colframe=teal!35!black,title={Before you move on}]
What does el granizo mean? (\textquotedblleft{}The hail\textquotedblright{}.) Say it once more.
\end{tcolorbox}
\section[granizar]{granizar — to hail}

\label{lesson:ES-C699-granizar}



\begin{quote}
Before the new one: say the Spanish for to flash with lightning, then the Spanish for the hail.
\end{quote}

\subsection*{You'll want to know: granizar}
\textbf{granizar} — \textquotedblleft{}to hail\textquotedblright{}.

\begin{quote}
Ayer granizó en la sierra.
\end{quote}

\begin{quote}
\emph{It hailed in the mountains yesterday.}
\end{quote}

Used in the third person, like \emph{llover}.

\begin{cousinweb}
\emph{granizar} is built on \emph{granizo} \textquotedblleft{}hail\textquotedblright{}, from \emph{grano} \textquotedblleft{}grain\textquotedblright{}, inherited from Latin \emph{grānum}.
\end{cousinweb}

\begin{grammarlens}[title={What you've built}]
One more word for telling, explaining and describing at B1.
\end{grammarlens}

\subsection*{Your turn}
\begin{itemize}
\raggedright
  \item \emph{Say it:} \emph{granizar}
  \item \emph{Say it:} \emph{granizar}, once more
  \item \emph{Say it:} a sentence of your own with \emph{granizar}
  \item \emph{Recall:} say the Spanish for to flash with lightning, then the Spanish for the hail, then say \emph{granizar} again
\end{itemize}

\begin{tcolorbox}[breakable,colback=teal!4,colframe=teal!35!black,title={Before you move on}]
What does granizar mean? (\textquotedblleft{}To hail\textquotedblright{}.) Say it once more.
\end{tcolorbox}
\section[tronar]{tronar — to thunder}

\label{lesson:ES-C699-tronar}



\begin{quote}
Before the new one: say the Spanish for the hail, then the Spanish for to hail.
\end{quote}

\subsection*{You'll want to know: tronar}
\textbf{tronar} — \textquotedblleft{}to thunder\textquotedblright{}.

\begin{quote}
Truena mucho esta noche.
\end{quote}

\begin{quote}
\emph{It is thundering a lot tonight.}
\end{quote}

The o becomes ue: \emph{truena}. \emph{El trueno} is thunder.

\begin{cousinweb}
\emph{tronar} is inherited from Latin \emph{tonāre} \textquotedblleft{}to thunder\textquotedblright{} (EtymDB), like \emph{trueno} \textquotedblleft{}thunder\textquotedblright{}.
\end{cousinweb}

\begin{grammarlens}[title={What you've built}]
One more word for telling, explaining and describing at B1.
\end{grammarlens}

\subsection*{Your turn}
\begin{itemize}
\raggedright
  \item \emph{Say it:} \emph{tronar}
  \item \emph{Say it:} \emph{tronar}, once more
  \item \emph{Say it:} a sentence of your own with \emph{tronar}
  \item \emph{Recall:} say the Spanish for the hail, then the Spanish for to hail, then say \emph{tronar} again
\end{itemize}

\begin{tcolorbox}[breakable,colback=teal!4,colframe=teal!35!black,title={Before you move on}]
What does tronar mean? (\textquotedblleft{}To thunder\textquotedblright{}.) Say it once more.
\end{tcolorbox}
